% ===================================================================
\section{Full-matrix training: a cross-entropy plateau}
\label{sec:fullmatrix}
\Cref{thm:separation} trains four tied scalars. Could a network whose every matrix entry trains escape the cancellation by learning other features,
other value maps or other attention patterns? Near the structured initialization it cannot. The proof needs no invariant manifold:
it combines the distributional cancellation of \cref{lem:matched} with a loss-floor and energy argument.

\paragraph{The full-matrix network.} Keep the embeddings and type mask of \cref{sec:architecture} fixed, and make every remaining entry trainable:
query and key $Q,K\in\mathbb R^{d\times7}$, value $V\in\mathbb R^{d_v\times7}$, output $O\in\mathbb R^{7\times d_v}$, MLP weights $U\in\mathbb R^{m\times7}$, bias $\bm a\in\mathbb R^m$,
and $D\in\mathbb R^{7\times m}$, and readout $\bm c\in\mathbb R^7$, with $d,d_v\ge1$ and $m\ge4$. The logit is
\[
 z_W=\bm c^\top\bigl[\bm g_b+D\relu(U\bm g_b+\bm a)\bigr],\qquad
 \bm g_b=\bm h_q^{(0)}+O\sum_{j=1}^N\alpha_jV\bm h_j^{(0)},\qquad
 \bm\alpha=\softmax\bigl((Q\bm h_q^{(0)})^\top K\bm h_j^{(0)}\bigr)_j .
\]
Let $W_0$ be the matrices of \cref{eq:attention-matrices,eq:mlp-matrices} at $\bm\theta(0)$, padded with zeros, with readout $\bm c=\bm e_7$; nothing is tied afterwards.
Training follows $\dot W\in-\partial\mathcal L_{\mathrm{ce}}(W)$, the Clarke subgradient flow, which is ordinary gradient flow wherever the loss is differentiable.

\begin{mainresult}
\begin{theorem}[Full-matrix cross-entropy plateau]
\label{thm:fullmatrix}
Under $\Dlaw0$, let $B\ge1$, $C_B=110(B+2)$ and $0<\varepsilon\le\tfrac{1}{10}$.
\begin{enumerate}[label=(\roman*),leftmargin=*,itemsep=2pt,topsep=2pt]
\item \emph{Weak credit near initialization.} Every $W$ with $\|W-W_0\|_F\le\varepsilon$ has
$|\E[Yz_W]|\le\gamma C_B\varepsilon^3$ and therefore $\mathcal L_{\mathrm{ce}}(W)\ge\log2-\tfrac12\gamma C_B\varepsilon^3$.
\item \emph{Plateau.} From $W(0)=W_0$, for all $0\le t<\tfrac{1}{C_B\varepsilon}$,
\[
 \|W(t)-W_0\|_F<\varepsilon,\qquad \mathcal L_{\mathrm{ce}}(W(t))\ge\log2-\tfrac12\gamma C_B\varepsilon^3 .
\]
Reaching any loss below $\log2-\delta$ with $\delta>\tfrac12\gamma C_B\varepsilon^3$ therefore takes time at least $\tfrac{1}{C_B\varepsilon}$.
\item \emph{Perturbed starts.} There is $L_B$ depending only on $B$ such that the same holds, with exit time at least $\tfrac{1}{4(C_B+(1+L_B)^2)\varepsilon}$,
from every $W(0)$ with $\|W(0)-W_0\|_F\le\varepsilon^{3/2}$.
\end{enumerate}
\end{theorem}
\end{mainresult}

\paragraph{Why the plateau holds.} At uniform attention the network sees only the instruction and the average embedding, which under $\Dlaw0$ is independent of the source, so
its logit has zero correlation with the label, whatever its other weights are (\cref{lem:matched}). Near $W_0$ attention can depart from uniform only by $O(\varepsilon)$,
and its dependence on the instruction, which the label needs, only by $O(\varepsilon^2)$: the key starts at zero and the query does not read the instruction.
Every path from attention to the logit passes through a product that starts at zero or at $O(\varepsilon)$, giving $|\E[Yz_W]|=O(\varepsilon^3)$.
Cross-entropy equals $\log2-\tfrac12\E[Yz]+\E\log\cosh\tfrac{z}{2}$, so the loss cannot fall by more than $O(\varepsilon^3)$ inside the ball, and the energy identity
$\|W(t)-W_0\|_F^2\le t\,[\mathcal L(W_0)-\mathcal L(W(t))]$ turns that into an exit time of order $\tfrac{1}{\varepsilon}$ (\cref{app:fullmatrix}).

\paragraph{What it does not say.} The theorem is a lower bound for the cancelled law only. It does not give the exact $\Theta(\tfrac{N}{\varepsilon})$ routing time,
the conserved quantity or the failure barrier for full matrices, and it does not prove that the independent law escapes fast; the
experiments in \cref{sec:numerics} compare the two. It covers the structured start and perturbations of size $\varepsilon^{3/2}$, not perturbations of fixed size
or fully random initialization. A loss near $\log2$ also does not by itself mean chance accuracy.

